\section{仿真}

本节为前述讨论提供数值证据。
仿真目的并非从数值上声称完整的安全定理，而是分离两个容易混淆的现象：
围捕误差的收敛性与无碰集合的前向不变性。
因此，我们对每个控制器在两种情形下进行测试：
第一种是良性情形，具有充分的初始间距、适中的初始速度、较强的 APF 排斥力
以及快速的滤波器；
第二种是压力情形，智能体彼此迎面运动、排斥力较弱、滤波器较慢。
压力情形旨在展示控制器在暂态后仍可能收敛，但此时已违反安全距离。

\subsection{APF 滤波滑模设计}

第一个仿真对应 APF 滤波滑模设计。
使用四个异构阶的智能体：
\[
    (m_1,m_2,m_3,m_4)=(1,2,3,4)
\]
良性情形采用标称参数
\[
    d=1.25,\quad \mu=10,\quad k_r=50,
    \quad \gamma_1=\cdots=\gamma_4=2.5,
\]
以及相对较大的滑模增益。
压力情形保持相同的智能体阶和安全距离，但降低滤波带宽和反馈强度，
并将第三阶和第四阶智能体初始化为近乎迎面的相对运动。

\begin{figure*}[t]
    \centering
    \begin{subfigure}[t]{0.24\textwidth}
        \centering
        \includegraphics[width=\linewidth]{../simulation/simulation3/out/safe_case/weighted_center_distance.pdf}
        \caption{仿真 3，良性情形}
        \label{fig:sim3-safe-center}
    \end{subfigure}
    \begin{subfigure}[t]{0.24\textwidth}
        \centering
        \includegraphics[width=\linewidth]{../simulation/simulation3/out/stress_case/weighted_center_distance.pdf}
        \caption{仿真 3，压力情形}
        \label{fig:sim3-stress-center}
    \end{subfigure}
    \begin{subfigure}[t]{0.24\textwidth}
        \centering
        \includegraphics[width=\linewidth]{../simulation/simulation4/out/safe_case/weighted_center_distance.pdf}
        \caption{仿真 4，良性情形}
        \label{fig:sim4-safe-center}
    \end{subfigure}
    \begin{subfigure}[t]{0.24\textwidth}
        \centering
        \includegraphics[width=\linewidth]{../simulation/simulation4/out/stress_case/weighted_center_distance.pdf}
        \caption{仿真 4，压力情形}
        \label{fig:sim4-stress-center}
    \end{subfigure}
    \caption{智能体加权中心到目标的距离。}
    \label{fig:simulation-center-distance}
\end{figure*}

\begin{figure*}[t]
    \centering
    \begin{subfigure}[t]{0.24\textwidth}
        \centering
        \includegraphics[width=\linewidth]{../simulation/simulation3/out/safe_case/pairwise_distances.pdf}
        \caption{仿真 3，良性情形}
        \label{fig:sim3-safe-pairwise}
    \end{subfigure}
    \begin{subfigure}[t]{0.24\textwidth}
        \centering
        \includegraphics[width=\linewidth]{../simulation/simulation3/out/stress_case/pairwise_distances.pdf}
        \caption{仿真 3，压力情形}
        \label{fig:sim3-stress-pairwise}
    \end{subfigure}
    \begin{subfigure}[t]{0.24\textwidth}
        \centering
        \includegraphics[width=\linewidth]{../simulation/simulation4/out/safe_case/pairwise_distances.pdf}
        \caption{仿真 4，良性情形}
        \label{fig:sim4-safe-pairwise}
    \end{subfigure}
    \begin{subfigure}[t]{0.24\textwidth}
        \centering
        \includegraphics[width=\linewidth]{../simulation/simulation4/out/stress_case/pairwise_distances.pdf}
        \caption{仿真 4，压力情形}
        \label{fig:sim4-stress-pairwise}
    \end{subfigure}
    \caption{智能体之间的成对距离。虚线为规定的安全距离。}
    \label{fig:simulation-pairwise-distance}
\end{figure*}

结果如
\cref{fig:simulation-center-distance,fig:simulation-pairwise-distance} 所示。
在良性情形下，加权中心距离收敛到零，且所有成对距离保持在安全距离之上。
最小智能体间距离为
\[
    \min_{i<j,t}\|p_i(t)-p_j(t)\|=1.593820>d=1.25 .
\]
因此，在此良性初始条件和足够响应的 APF 滤波器下，仿真同时满足围捕
要求和数值避碰要求。

然而，在压力情形下，同一控制器不再能保持安全性。
加权中心误差最终收敛到较小值，但最小成对距离满足
\[
    \min_{i<j,t}\|p_i(t)-p_j(t)\|=0.000975<d=1.25 ,
\]
由智能体 $3$ 和 $4$ 在 $t\approx0.64$ 秒附近达到。
因此，滤波滑模变量的收敛，或暂态后加权中心的收敛，并不意味着无碰集合
是前向不变的。
这与理论分析中的障碍完全相同：APF 项在到达高阶输入通道之前被滤波，
因此当智能体以大相对速度或高阶惯性趋近边界时，奇异排斥效应可能来得太晚。

\subsection{指令滤波反步设计}

第二个仿真对应指令滤波反步尝试。
使用五个智能体，其阶为
\[
    (m_1,m_2,m_3,m_4,m_5)=(3,3,3,4,4).
\]
良性情形使用较小的滤波时间常数和适中的初始速度。
压力情形使用较慢的指令滤波器、较弱的排斥力以及近乎迎面的初始条件。
该情形直接测试结果 4 中识别出的扰动项
\[
    r_i=e_{1,i}+\varepsilon_{1,i}
\]

对于良性情形，最小成对距离为
\[
    \min_{i<j,t}\|p_i(t)-p_j(t)\|=3.525914>d=1 .
\]
加权中心误差收敛到零，且目标在所有采样时刻保持在凸包内。
这表明当理想的阶一 APF 运动被足够精确地跟踪时，
即当残余 $r_i$ 保持较小且初始间距给控制器足够的反应时间时，
指令滤波反步设计可以良好工作。

对于具有 $100$ 秒时间 horizon 的压力情形，行为则非常不同。
最小成对距离为
\[
    \min_{i<j,t}\|p_i(t)-p_j(t)\|=0.011129<d=1 ,
\]
在智能体 $1$ 和 $2$ 之间、约 $t=95.12$ 秒附近达到。
在此运行中残余变得非常大，且沿实际滤波闭环的理想 APF 能量增加而不是
保持单调。
这在数值上吻合结果 4 的估计，
\[
    \dot W
    =
    -\sum_i\|\kappa_i p_{i0}-\phi_i\|^2
    +\sum_i(\kappa_i p_{i0}-\phi_i)^T r_i ,
\]
其中第二项可以主导标称耗散。

\subsection{讨论}

这些仿真表明，仅靠跟踪变量的收敛性之外，还需要额外条件才能保证安全性。
具体而言，良性情形具有以下共同特征：
\begin{itemize}
    \item 初始构型与碰撞边界有明显的距离裕量；
    \item 初始相对速度和高阶状态并非激进地指向边界；
    \item APF 或指令滤波器足够快，使得排斥项不会被过度延迟；
    \item 反馈增益足够大，使得暂态期间实现的阶一 APF 残余较小；
    \item 数值屏障相对于趋近速度足够强。
\end{itemize}
在此条件下，轨迹保持在安全集合的紧致子集内，因此 APF 导数和滤波命令
是正则的。
闭环行为接近理想的阶一 APF 运动，且数值上观察到避碰。

压力情形恰恰违反了这些非正式条件。
智能体可以以高于高阶闭环将奇异 APF 信号转化为位置级分离的速度
趋近安全边界。
因此，控制器在暂态后仍可能恢复收敛，但安全性已经丧失。
因此，这些仿真支持本笔记的主要结论：基于 APF 的高阶跟踪控制器在有利
情形下可以满足围捕，但严格的避碰保证需要额外的前向不变性机制，
例如直接控制碰撞边界的高阶 CBF 层或障碍李雅普诺夫设计。